% formal/axiomatic.tex
% mainfile: ../perfbook.tex
% SPDX-License-Identifier: CC-BY-SA-3.0

\section{公理化方法}
\label{sec:formal:Axiomatic Approaches}
\OriginallyPublished{sec:formal:Axiomatic Approaches}{Axiomatic Approaches}{Linux Weekly News}{PaulEMcKenney2014weakaxiom}
%
\epigraph{理论帮助我们承受对事实的无知。}
	{George Santayana}

虽然 PPCMEM 工具可以求解如
\cref{lst:formal:IRIW Litmus Test}所示的著名
``independent reads of independent writes''（IRIW）litmus test，
但这样做需要不少于14个 CPU 小时，
并产生不少于10 GB 的 state space。
话虽如此，这种情况相比 PPCMEM 出现之前已经大有改善，
在那之前解决这个问题需要翻阅大量参考手册、
尝试证明、与专家讨论，
然后在所有这些之后，仍然怀疑最终答案。
虽然14小时看起来可能很长，但它比几周甚至几个月
要短得多。

\begin{listing}
\begin{fcvlabel}[ln:formal:IRIW Litmus Test]
\begin{VerbatimL}[commandchars=\%\@\$]
PPC IRIW.litmus
""
(* Traditional IRIW. *)
{
0:r1=1; 0:r2=x;
1:r1=1;         1:r4=y;
        2:r2=x; 2:r4=y;
        3:r2=x; 3:r4=y;
}
P0           | P1           | P2           | P3           ;
stw r1,0(r2) | stw r1,0(r4) | lwz r3,0(r2) | lwz r3,0(r4) ;
             |              | sync         | sync         ;
             |              | lwz r5,0(r4) | lwz r5,0(r2) ;

exists
(2:r3=1 /\ 2:r5=0 /\ 3:r3=1 /\ 3:r5=0)
\end{VerbatimL}
\end{fcvlabel}
\caption{IRIW Litmus Test}
\label{lst:formal:IRIW Litmus Test}
\end{listing}

然而，鉴于 litmus test 的简单性，所需的时间和空间
有点令人惊讶——它有两个线程存储到两个不同的变量，
另外两个线程以相反的顺序从这两个变量加载。
如果两个加载线程对两个存储的顺序有分歧，
则 assertion 触发。
即使按照 memory-order litmus test 的标准，这也是相当简单的。

PPCMEM 消耗大量时间和空间的一个原因是它进行了
基于轨迹的完整 state-space search，这意味着它必须生成
和评估在架构级别上事件和动作的所有可能顺序和组合。
在这个级别上，加载和存储都对应于精心设计的
事件和动作序列，导致必须完全搜索的非常大的 state space，
进而导致大量的内存和 CPU 消耗。

当然，许多轨迹彼此非常相似，这表明
将相似轨迹视为一个的方法可能会
提高性能。
一种这样的方法是 \pplsur{Jade}{Alglave}
等人~\cite{Alglave:2014:HCM:2594291.2594347}
的 axiomatic approach，它创建一组公理来表示 \IX{memory model}，
然后将 litmus test 转换为可以在这组公理上
证明或反证的定理。
由此产生的工具称为 \qco{herd}，方便地接受与 PPCMEM
相同的 litmus test 作为输入，包括
\cref{lst:formal:IRIW Litmus Test}中所示的 IRIW litmus test。

\begin{listing}
\begin{fcvlabel}[ln:formal:Expanded IRIW Litmus Test]
\begin{VerbatimL}[commandchars=\%\@\$]
PPC IRIW5.litmus
""
(* Traditional IRIW, but with five stores instead of *)
(* just one.                                         *)
{
0:r1=1; 0:r2=x;
1:r1=1;         1:r4=y;
        2:r2=x; 2:r4=y;
        3:r2=x; 3:r4=y;
}
P0           | P1           | P2           | P3           ;
stw r1,0(r2) | stw r1,0(r4) | lwz r3,0(r2) | lwz r3,0(r4) ;
addi r1,r1,1 | addi r1,r1,1 | sync         | sync         ;
stw r1,0(r2) | stw r1,0(r4) | lwz r5,0(r4) | lwz r5,0(r2) ;
addi r1,r1,1 | addi r1,r1,1 |              |              ;
stw r1,0(r2) | stw r1,0(r4) |              |              ;
addi r1,r1,1 | addi r1,r1,1 |              |              ;
stw r1,0(r2) | stw r1,0(r4) |              |              ;
addi r1,r1,1 | addi r1,r1,1 |              |              ;
stw r1,0(r2) | stw r1,0(r4) |              |              ;

exists
(2:r3=1 /\ 2:r5=0 /\ 3:r3=1 /\ 3:r5=0)
\end{VerbatimL}
\end{fcvlabel}
\caption{扩展的 IRIW Litmus Test}
\label{lst:formal:Expanded IRIW Litmus Test}
\end{listing}

然而，PPCMEM 需要14个 CPU 小时来求解 IRIW，\co{herd} 只需
17毫秒即可完成，这代表了超过六个数量级的加速。
话虽如此，问题本质上是指数级的，所以我们应该预期
\co{herd} 对更大的问题会表现出指数级的减慢。
这正是所发生的，例如，如果我们为每个写入 CPU 再添加四次写入，
如\cref{lst:formal:Expanded IRIW Litmus Test}所示，
\co{herd} 减慢了超过50,000倍，需要超过
15\emph{分钟}的 CPU 时间。
添加线程也会导致指数级的
减慢~\cite{PaulEMcKenney2014weakaxiom}。

尽管具有指数级的性质，PPCMEM 和 \co{herd} 都已被证明
对检查关键的并行算法非常有用，包括 x86 系统上的
队列锁交接。
\co{herd} 工具的弱点与 PPCMEM 类似，
后者在\cref{sec:formal:PPCMEM Discussion}中描述。
有一些隐蔽（但非常真实）的情况，PPCMEM 和
\co{herd} 工具不一致，截至2021年，其中许多但不是全部分歧
已被解决。

如果 litmus test 可以用 C 语言编写
（如\cref{lst:formal:Meaning of PPCMEM Litmus Test}）
而不是汇编语言
（如\cref{lst:formal:PPCMEM Litmus Test}），
那将会很有帮助。
这现在是可能的，将在以下各节中描述。

\subsection{Axiomatic Approaches 和锁}
\label{sec:formal:Axiomatic Approaches and Locking}

Axiomatic approaches 也可以应用于更高级别的语言
以及更高级别的同步原语，
如\cref{lst:formal:Locking Example}（\path{C-Lock1.litmus}）
中所示的基于锁的 litmus test 所例证的。
这个 litmus test 可以由
\IXalthmr{\acrf{lkmm}}{Linux kernel}{memory model}~%
\cite{Alglave:2018:FSC:3173162.3177156,LucMaranget2018lock.cat}建模。
如预期的那样，\co{herd} 工具的输出包含字符串 \co{Never}，
正确地表示 \co{P1()} 不能看到 \co{x} 具有值
1。\footnote{
	\co{herd} 工具的输出与 PPCMEM 兼容，
	所以可以查看
	\cref{lst:formal:PPCMEM Detects an Error,%
	lst:formal:PPCMEM on Repaired Litmus Test}
	中展示输出格式的示例。}

\begin{listing}
\input{CodeSamples/formal/herd/C-Lock1@whole.fcv}
\caption{锁示例}
\label{lst:formal:Locking Example}
\end{listing}

\QuickQuiz{
	你需要做什么才能在如
	\cref{lst:formal:Locking Example}所示的 litmus test 上运行 \co{herd}？
}\QuickQuizAnswer{
	获取 Linux 内核源代码的 v4.17（或更高）版本，
	然后按照 \path{tools/memory-model/README} 中的说明
	安装所需的工具。
	然后按照进一步的说明在你选择的 litmus test 上
	运行这些工具。
}\QuickQuizEnd

\begin{listing}
\input{CodeSamples/formal/herd/C-Lock2@whole.fcv}
\caption{有缺陷的锁示例}
\label{lst:formal:Broken Locking Example}
\end{listing}

当然，如果 \co{P0()} 和 \co{P1()} 使用不同的锁，如
\cref{lst:formal:Broken Locking Example}（\path{C-Lock2.litmus}）
所示，那么一切都不确定了。
在这种情况下，\co{herd} 工具的输出包含字符串
\co{Sometimes}，正确地表示使用不同的锁允许
\co{P1()} 看到 \co{x} 具有值1。

\QuickQuiz{
	为什么要费心直接建模锁？
	为什么不简单地用原子操作模拟锁？
}\QuickQuizAnswer{
	一个词，性能，如
	\cref{tab:formal:Locking: Modeling vs. Emulation Time (s)}中所示。
	第一列显示建模的 \co{herd} 进程数量。
	第二列显示直接在 \co{herd} 的 cat 语言中建模
	\co{spin_lock()} 和 \co{spin_unlock()} 时的 \co{herd} 运行时间。
	第三列显示使用 \co{cmpxchg_acquire()} 模拟 \co{spin_lock()}
	和使用 \co{smp_store_release()} 模拟 \co{spin_unlock()} 时的
	\co{herd} 运行时间，使用 \co{herd} 的 \co{filter} 子句
	拒绝未能获取锁的执行。
	第四列与第三列类似，但使用 \co{xchg_acquire()} 而不是
	\co{cmpxchg_acquire()}。
	第五列和第六列与第三列和第四列类似，
	但改用 \co{herd} 的 \co{exists} 子句来拒绝
	未能获取锁的执行。

\begin{table}
\rowcolors{10}{}{lightgray}
\small
\centering
\newcommand{\lockfml}[1]{\multicolumn{1}{c}{\begin{picture}(6,8)(0,0)\rotatebox{90}{#1}\end{picture}}}
\begin{tabular}{rrrrrr}
	\toprule
	& 建模 & \multicolumn{4}{c}{模拟} \\
	\cmidrule(l){2-2} \cmidrule(l){3-6}
	& & \multicolumn{2}{c}{\tco{filter}} & \multicolumn{2}{c}{\tco{exists}} \\
	\cmidrule(l){3-4} \cmidrule(l){5-6}
	\lockfml{进程数}
	&
	& \tco{cmpxchg}
	& \tco{xchg}
	& \tco{cmpxchg}
	& \tco{xchg}
	\\
	\cmidrule{1-1} \cmidrule(l){2-2} \cmidrule(l){3-4} \cmidrule(l){5-6}
	2 & 0.004 &  0.022 &  0.027 &   0.039 &   0.058 \\
	3 & 0.041 &  0.743 &  0.968 &   1.653 &   3.203 \\
	4 & 0.374 & 59.565 & 74.818 & 151.962 & 500.960 \\
	5 & 4.905 \\
	\bottomrule
\end{tabular}
\caption{锁：建模 vs.\@ 模拟时间（秒）}
\label{tab:formal:Locking: Modeling vs. Emulation Time (s)}
\end{table}

	还要注意，使用 \co{filter} 子句大约比使用
	\co{exists} 子句快两倍。
	这并不奇怪，因为 \co{filter} 子句允许
	早期放弃被排除的执行，而被排除的执行
	是锁被多个进程同时持有的执行。

	更重要的是，直接建模 \co{spin_lock()} 和 \co{spin_unlock()}
	比建模模拟的锁快5倍到超过两个数量级。
	这也不应该令人惊讶，因为直接建模提高了
	抽象级别，从而减少了 \co{herd} 必须建模的事件数量。
	因为 \co{herd} 所做的几乎所有事情都具有指数级的
	计算复杂度，事件数量的适度减少
	会产生指数级的运行时间缩减。

	因此，在 formal verification 中甚至比在并行编程本身中
	更要强调分而治之！！！
}\QuickQuizEnd

但锁不是唯一可以直接建模的同步原语：
下一节将介绍 RCU\@。

\subsection{Axiomatic Approaches 和 RCU}
\label{sec:formal:Axiomatic Approaches and RCU}

\begin{fcvref}[ln:formal:C-RCU-remove:whole]
Axiomatic approaches 也可以分析涉及
RCU 的 litmus test~\cite{Alglave:2018:FSC:3173162.3177156}。
为此，\cref{lst:formal:Canonical RCU Removal Litmus Test}
（\path{C-RCU-remove.litmus}）
展示了一个对应于经典的 RCU 介导的从链表中删除操作的
litmus test。
与锁 litmus test 一样，这个 RCU litmus test 可以
由 LKMM 建模，与建模 RCU 的模拟相比具有
类似的性能优势。
\Clnref{head}将 \co{x} 显示为列表头，初始
引用 \co{y}，后者在\clnref{tail:1}上被初始化为值 \co{2}。

\begin{listing}
\input{CodeSamples/formal/herd/C-RCU-remove@whole.fcv}
\caption{经典 RCU 删除 Litmus Test}
\label{lst:formal:Canonical RCU Removal Litmus Test}
\end{listing}

\Clnrefrange{P0start}{P0end}上的 \co{P0()}
通过将元素 \co{y} 替换为元素 \co{z}
（\clnref{assignnewtail}）来从列表中删除元素 \co{y}，
等待一个 \IX{grace period}（\clnref{sync}），
最后将 \co{y} 清零以模拟 \co{free()}（\clnref{free}）。
\Clnrefrange{P1start}{P1end}上的 \co{P1()}
在一个 RCU 读端临界区内执行
（\clnrefrange{rl}{rul}），
获取列表头（\clnref{rderef}），然后
加载下一个元素（\clnref{read}）。
下一个元素应该是非零的，即尚未被释放
（\clnref{exists_}）。
几个其他变量被输出用于调试目的
（\clnref{locations_}）。

\co{herd} 工具运行此 litmus test 时的输出包含
\co{Never}，表示 \co{P0()} 从不访问已释放的元素，
这是预期的。
同样如预期的那样，删除\clnref{sync}会导致 \co{P0()}
访问已释放的元素，如 \co{herd} 输出中的 \co{Sometimes}
所指出的。
\end{fcvref}

\begin{listing}
\ebresizeverb{.98}{\input{CodeSamples/formal/herd/C-RomanPenyaev-list-rcu-rr@whole.fcv}}
\caption{复杂的 RCU Litmus Test}
\label{lst:formal:Complex RCU Litmus Test}
\end{listing}

\begin{fcvref}[ln:formal:C-RomanPenyaev-list-rcu-rr:whole]
\ppl{Roman}{Penyaev}~\cite{RomanPenyaev2018rrRCU}提出的
一个更复杂的例子的 litmus test 如
\cref{lst:formal:Complex RCU Litmus Test}
（\path{C-RomanPenyaev-list-rcu-rr.litmus}）所示。
在这个例子中，读者（由\clnrefrange{P0start}{P0end}上的
\co{P0()} 建模）以轮询方式访问链表，
将指向最后访问的列表元素的指针
``泄漏''到变量 \co{c} 中。
更新者（由\clnrefrange{P1start}{P1end}上的 \co{P1()} 建模）
删除一个元素，注意不要干扰当前或未来的读者。

\QuickQuiz{
	等等！！！
	从 RCU 读端临界区泄漏指针
	不是一个严重的 bug 吗？？？
}\QuickQuizAnswer{
	是的，它通常是一个严重的 bug。
	然而，在这种情况下，更新者被巧妙地构造为
	正确处理此类指针泄漏。
	但请不要养成做这种事情的习惯，
	特别是在没有花大量心思使某种更常规的方法
	可行之前不要这样做。
}\QuickQuizEnd

\Clnrefrange{listtail}{listhead}定义了初始链表，
先从尾部开始。
在 Linux 内核中，这将是一个双向链接的循环列表，
但 \co{herd} 目前无法建模这样的结构。
因此策略是使用一个足够长的单链接线性列表，
使得永远不会到达末尾。
\Clnref{rrcache}定义了变量 \co{c}，用于在连续的
RCU 读端临界区之间缓存列表指针。

同样，\clnrefrange{P0start}{P0end}上的 \co{P0()} 建模读者。
这个进程建模了一对连续的读者以轮询方式遍历列表，
第一个读者在\clnrefrange{rl1}{rul1}上，
第二个读者在\clnrefrange{rl2}{rul2}上。
\Clnref{rdcache}获取缓存在 \co{c} 中的指针，
如果\clnref{rdckcache}看到指针为 \co{NULL}，
\clnref{rdinitcache}从列表开头重新开始。
无论哪种情况，\clnref{rdnext}前进到下一个列表元素，
\clnref{rdupdcache}将指向该元素的指针存储回
变量 \co{c} 中。
\Clnrefrange{rl2}{rul2}重复此过程，但使用
寄存器 \co{r3} 和 \co{r4} 而不是 \co{r1} 和 \co{r2}。
与\cref{lst:formal:Canonical RCU Removal Litmus Test}一样，
这个 litmus test 存储零来模拟 \co{free()}，所以
\clnref{exists_}检查这四个寄存器中是否有任何一个为
\co{NULL}，也就是零。

因为 \co{P0()} 从其第一个 RCU 读端临界区
向第二个泄漏了一个受 RCU 保护的指针，\co{P1()} 必须
非常小心地执行其更新（删除 \co{x}）。
\Clnref{updremove}通过将 \co{w} 链接到 \co{y} 来删除 \co{x}。
\Clnref{updsync1}等待读者，之后没有后续读者
可以通过链表找到通往 \co{x} 的路径。
\Clnref{updrdcache}获取 \co{c}，如果\clnref{updckcache}
确定 \co{c} 引用了新删除的 \co{x}，
\clnref{updinitcache}将 \co{c} 设置为 \co{NULL}，
\clnref{updsync2}再次等待读者，之后
后续读者无法从 \co{c} 获取 \co{x}。
无论哪种情况，\clnref{updfree}通过存储零到 \co{x}
来模拟 \co{free()}。

\QuickQuiz{
	\begin{fcvref}[ln:formal:C-RomanPenyaev-list-rcu-rr:whole]
	在\cref{lst:formal:Complex RCU Litmus Test}中，
	为什么读者不能在 \co{P1()} 在\clnref{updinitcache}
	上将 \co{c} 清零之前获取它，然后
	在它被清零之后将相同的值存储回 \co{c}，
	从而击败清零操作？
	\end{fcvref}
}\QuickQuizAnswer{
	\begin{fcvref}[ln:formal:C-RomanPenyaev-list-rcu-rr:whole]
	因为读者在\clnref{rdnext}上前进到下一个元素，
	从而避免存储指向与获取的元素相同的指针。
	\end{fcvref}
}\QuickQuizEnd

\co{herd} 工具运行此 litmus test 时的输出包含
\co{Never}，表示 \co{P0()} 从不访问已释放的元素，
这是预期的。
同样如预期的那样，删除任一 \co{synchronize_rcu()}
会导致 \co{P1()} 访问已释放的元素，如 \co{herd}
输出中的 \co{Sometimes} 所指出的。
\end{fcvref}

\QuickQuizSeries{%
\QuickQuizB{
	\begin{fcvref}[ln:formal:C-RomanPenyaev-list-rcu-rr:whole]
	在\cref{lst:formal:Complex RCU Litmus Test}中，
	为什么不在\clnref{updfree}之前只调用一次
	\co{synchronize_rcu()}？
	\end{fcvref}
}\QuickQuizAnswerB{
	\begin{fcvref}[ln:formal:C-RomanPenyaev-list-rcu-rr:whole]
	因为这会导致 \co{P0()} 访问已释放的元素。
	但不要只听我说，在 \co{herd} 中试一试！
	\end{fcvref}
}\QuickQuizEndB
%
\QuickQuizE{
	\begin{fcvref}[ln:formal:C-RomanPenyaev-list-rcu-rr:whole]
	同样在\cref{lst:formal:Complex RCU Litmus Test}中，
	\clnref{updfree}不能用 \co{WRITE_ONCE()} 而不是
	\co{smp_store_release()} 吗？
	\end{fcvref}
}\QuickQuizAnswerE{
	\begin{fcvref}[ln:formal:C-RomanPenyaev-list-rcu-rr:whole]
	这是一个很好的问题。
	截至2021年末，答案是``没人知道''。
	这在很大程度上取决于 \ARMv8 的条件移动指令的语义。
	在等待这些语义明确之前，\co{smp_store_release()}
	是安全的选择。
	\end{fcvref}
}\QuickQuizEndE
}

这些章节展示了 axiomatic approaches 如何能够成功地
在 C 语言 litmus test 中建模锁和 RCU 等同步原语。
从更长远来看，希望 axiomatic approaches 能建模
更高级别的软件工件，产生指数级的
验证加速。
这可能允许对更大的软件系统进行公理化验证，
也许结合分离逻辑中的空间同步
技术~\cite{AlexeyGotsman2013ESOPRCU,PeterWOHearn2001SeparationLogic}。
另一种选择是将布尔逻辑的公理投入使用，
如下一节所描述的。
